\section{Discussion}
\label{sec:discussion}

\subsection{Key Findings}

The DNSI-gap correlation ($R = -0.95$) suggests that improvement entropy, normalized by task difficulty, may capture benchmark evolution dynamics. However, with $n=4$, this represents preliminary evidence. The temporal prediction result ($R^2 = 0.35$) suggests potential as a leading indicator, though LOO-CV instability ($R^2 = -3.82$) indicates the model does not yet generalize reliably --- larger benchmark samples are needed before practical deployment.

\subsection{Limitations}

\textbf{Small Sample Size ($n = 4$):} Bootstrap CIs span the full range; we frame results as pilot study.

\textbf{Synthetic SOTA Histories:} Proof-of-concept uses historically-accurate synthetic data; production implementation should use real PapersWithCode data.

\textbf{NLP Difficulty Proxy Undefined:} Class count maps less naturally to NLP tasks; future work should develop domain-specific normalizers.

\textbf{Correlation $\neq$ Causation:} We demonstrate predictive correlation, not causal mechanism.

\subsection{Broader Impact}

DNSI could help researchers allocate effort more effectively by identifying saturated benchmarks. The metric could be gamed by artificially diversifying improvement patterns; we recommend DNSI as one input among many for benchmark assessment.
